\section{Program Implementation} \label{sec:implementation}

ENEF provided training on how to incorporate financial literacy textbooks into the school curriculum of the 2015 school year (that started in February and ended at the beginning of December). In Joinville, the training was held in February 2015 for the teaching supervisors of each school. In Manaus, supervisors of the regional education boards were trained at the end of March of the same year. Teachers in the schools that implemented the intervention were trained by those who attended the ENEF training. The pilot schools had the autonomy to decide when to introduce the textbooks into their classes. 

The programs' supervisors designed by the Departments of Education of Joinville and Manaus were in charge of monitoring the pilot implementation. During the 2015 school year, the supervisors filled out questionnaires aimed to identify the major bottlenecks and keep track of the use of financial literacy textbooks. ENEF used the findings of the first two questionnaires applied during the first semester of the school year to improve the program's implementation. At the beginning of the second semester, ENEF organized a meeting with teachers in charge of incorporating the pilot textbooks into their classes. This meeting sought to stimulate the use of the textbooks and to encourage teachers who were already using them to share their experiences with those who had had little contact with the material at that point.

Insights into the programs' implementation can also be obtained from the questionnaires applied to teachers and students at the end of the 2015 school year. \autoref{tab:implementation} shows that around 80\% of pilot students answered the financial literacy assessment. With the exception of third graders, the socioeconomic questionnaire was answered by at least 97\% of those who did the financial literacy test. We observe that there are no significant differences in response rates of treatment and control groups. 

\autoref{tab:implementation} shows that almost half of the third-grade teachers introduced financial literacy textbooks during the whole school year. However, middle school teachers only introduced the textbooks in the second semester. Also, according to questionnaires applied to students, more than half of middle-school students received financial literacy textbooks in the second semester.  Slightly more than 20\% of middle-school students had at least 80\% of the financial literacy content covered. The late introduction of the textbooks is reflected in the percentage of the curricula actually covered. One can infer from the data on the program's implementation that the treatment was not very intense, and its intensity differed across schools and grades.  
